% Reconstructed author proof on 2026-10-21 after workspace loss.
% This is newly written source, not a byte-identical recovery of the lost
% 54d3814d603f95ce3f2b1ec3238cf1bb05bfa47fd35218ec0ed56fb368ebb3cc.
% The frozen inference program and its prospective rules are unchanged.
\section{Joint inference for the fixed original policies}
\label{sec:inference122}

This section gives the probability argument implemented by
\texttt{code122/stats122.py}. It concerns costs of the original investment
problem and the exact fixed catalogue. It does not provide a population
adoption theorem for a neural representation. All statements condition on
the information available before the new state acquisition, including the
inherited actors, the frozen source, and its deterministic rules. They do
not condition on the realized values of the subsequently acquired states.

\paragraph{Population, numerical premises, and the economic random variables.}
Let
\[
 \mathcal{X}_{40}=
 \left\{0,2^{-40},\ldots,(2^{40}-1)2^{-40}\right\}^{2},
 \qquad P_{40}=\operatorname{Unif}(\mathcal{X}_{40}).
\]
Assume the prescribed entropy acquisition supplies independent uniform
40-bit coordinates, independently of the pre-acquisition information.
The sample is then \(X_1,\ldots,X_n\) independently distributed according
to \(P_{40}\), with \(n=32\). Sampling is with replacement; no state is
excluded or replaced because of its economic outcome. Both risk levels
and all methods use these same state units. The original economic
innovations remain continuously uniform on their declared interval.
Discretization of the initial-state sampling law does not replace that
continuous innovation law.

For a fixed risk index and state \(x\), let \(V(x)\) denote the optimal
value of the original problem and let \(J_m(x)\) be the original complete
policy cost delivered by method \(m\), including its prescribed future
policy. The statistical premises are that these policies are
well-defined measurable state functions and, uniformly on
\(\mathcal{X}_{40}\),
\begin{equation}
 V(x)\le J_m(x)\le V(x)+\tau,\qquad \tau=2^{-8}.
 \label{eq:inference-support122}
\end{equation}
Every policy uses the same original objective, information set, feasible
decisions, and interpretation of its future-policy cost. In particular,
agreement of first actions can be used as a cost identity only because
the compared methods also use the same prescribed future policy.

For each of the 58 actor/risk cells, denote its separately paid control
by \(m0\), and its recycled and adaptive services by \(me\).
The support-to-action safeguards prove, state by state,
\begin{equation}
 0\le c_{me}(x)\le
 G_{me}(x):=J_{m0}(x)-J_{me}(x)\le\tau.
 \label{eq:inference-safe-gain122}
\end{equation}
Here \(c_{me}\) is the reported accumulated certified gain; on retention
of the incumbent it is zero. The upper bound follows from
\eqref{eq:inference-support122}, not from an actor fitting loss.
For any two qualified methods \(m,b\) at the same risk index,
\begin{equation}
 H_{mb}(x):=J_b(x)-J_m(x)\in[-\tau,\tau].
 \label{eq:inference-signed-support122}
\end{equation}
Consequently the support for signed comparisons has twice the width of
the support for guaranteed own-control improvement.

The independent evaluator supplies valid outward enclosures of the
same original complete-policy costs and their paired differences.
Writing its raw safe-gain interval as
\([\underline g^{\,I}_{mj},\overline g^{\,I}_{mj}]\), the interval used by
the frozen statistics program is
\begin{equation}
 \ell_{mj}=
 \max\{0,c_m(X_j),\underline g^{\,I}_{mj}\},\qquad
 u_{mj}=\min\{\tau,\overline g^{\,I}_{mj}\}.
 \label{eq:inference-intersection122}
\end{equation}
Under the stated numerical premises,
\(\ell_{mj}\le G_m(X_j)\le u_{mj}\). The raw independent interval and the
proof-supported intersection are distinct objects and are both
reported. Taking their intersection does not independently validate the
proof-derived endpoint. A conflicting or missing enclosure is a failure
to satisfy the numerical premises; it cannot be repaired by deleting a
sample or by relabelling an unresolved independent comparison as a
confirmation.

For a signed comparison, subtract the independent policy enclosures in
the cost-difference direction and intersect with \([-\tau,\tau]\).
If the saved exact action identity and the common future-policy rule
prove equality on a sampled state, its paired interval is \([0,0]\).
This equality for finitely many sampled states does not make the
population difference of two distinct methods identically zero.
The frozen program uses the exact population interval \([0,0]\) only
when the full method key equals the comparator key.

\paragraph{Computational domain of the uniform premise.}
The all-state premises here concern the dyadic initial-state support
\(\mathcal{X}_{40}\). The finite exact-rational service argument in
Section~\ref{sec:proofs96}, together with the original feasible fallback
and policy qualification, is a mathematical totalization on that
domain. It is not an observed uniform wall-time or memory bound.
Infrastructure failures and incomplete attempts remain recorded; no
operational failure is silently converted to an extra accepted
observation. A transfer from \(P_{40}\) to a continuously uniform
initial-state population would require an independently justified
modulus for the actual implemented policy cost, or a separate acquisition
and rounding argument. Branching, acceptance, and stopping decisions can
make the implemented policy discontinuous, so continuity of the original
primitive alone does not supply this transfer. These qualifications do
not change the continuous economic innovation distribution.

\begin{lemma}[A bounded-mean inequality with interval observations]
\label{lem:bounded-mean122}
Let \(Y_1,\ldots,Y_n\) be independent copies of \(Y\in[a,b]\), \(a<b\),
and put \(\mu=\mathbb E Y\). For \(r>0\),
\[
 \mathbb P\{|\overline Y-\mu|>r\}
 \le 2\exp\{-2nr^2/(b-a)^2\}.
\]
If data-dependent outward endpoints obey
\(\ell_j\le Y_j\le u_j\) for every sample, then, with at least the same
coverage, \(\mu\) belongs to
\[
 [\,\max\{a,\overline\ell-r\},
       \min\{b,\overline u+r\}\,].
\]
No independence of the endpoints from their own observations is needed.
\end{lemma}

\begin{proof}
For any bounded random variable \(Z\) on \([a,b]\), under any probability
law on that interval,
\[
 \operatorname{Var}(Z)
 \le \mathbb E\bigl(Z-(a+b)/2\bigr)^2
 \le (b-a)^2/4.
\]
For real \(t\), set
\(\psi(t)=\log\mathbb E\exp(t(Y-\mu))\). Boundedness permits
differentiation under the expectation. The second derivative is the
variance of \(Y\) under the exponentially tilted probability law, hence
\(0\le\psi''(t)\le(b-a)^2/4\). Since
\(\psi(0)=\psi'(0)=0\), the integral remainder formula gives, also for
negative \(t\),
\[
 \psi(t)=t^2\int_0^1(1-s)\psi''(st)\,ds
 \le t^2(b-a)^2/8.
\]
Independence and Markov's inequality therefore imply, for \(\lambda>0\),
\[
 \mathbb P\{\overline Y-\mu\ge r\}
 \le \exp\{-\lambda nr+n\lambda^2(b-a)^2/8\}.
\]
Choosing \(\lambda=4r/(b-a)^2\) gives the upper-tail bound; replacing
\(Y\) by \(-Y\) gives the lower-tail bound. Their union gives the
displayed two-sided inequality (and also controls strict failure).
Finally \(\overline\ell\le\overline Y\le\overline u\). On the event
\(|\overline Y-\mu|\le r\), the mean lies between
\(\overline\ell-r\) and \(\overline u+r\); the known support justifies
clipping to \([a,b]\).
\end{proof}

\begin{lemma}[Exact one-sided event inversion]
\label{lem:binomial-inversion122}
Let \(K\sim\operatorname{Bin}(n,p)\), fix \(0<\delta<1\), and define
\(L_0=0\). For \(1\le k\le n\), let \(L_k\) be the unique root in
\((0,1)\) of
\[
 B_k(v):=\sum_{j=k}^n {n\choose j}v^j(1-v)^{n-j}=\delta.
\]
Then \(\mathbb P_p\{L_K>p\}\le\delta\). Replacing every \(L_k\) by a
smaller nonnegative number preserves the bound.
\end{lemma}

\begin{proof}
For \(1\le k\le n\), \(B_k(0)=0\), \(B_k(1)=1\), and differentiation
and cancellation give
\[
 B_k'(v)=n{n-1\choose k-1}v^{k-1}(1-v)^{n-k}>0
 \quad(0<v<1).
\]
Thus the root exists and is unique. Since \(B_{k+1}(v)\le B_k(v)\),
the sequence \(L_k\) is nondecreasing. For \(p\in(0,1)\), if there is
any \(k\) with \(L_k>p\), let \(k_0\) be the smallest such index.
Then \(\{L_K>p\}=\{K\ge k_0\}\), whose probability is
\(B_{k_0}(p)<\delta\). If no such \(k\) exists the probability is zero.
At \(p=0\), \(K=0\); at \(p=1\), no \(L_k\) exceeds \(p\).
Decreasing the endpoints only decreases the failure event.
\end{proof}

\paragraph{Rational implementation of the event endpoints.}
The frozen program uses \(\delta=1/23200\) and 48 bisections. For
\(k\ge1\), initialize \(l=0,h=1\). At a rational midpoint \(v=c/d\),
test the exact integer inequality
\begin{equation}
 23200\sum_{j=k}^{32}{32\choose j}c^j(d-c)^{32-j}
 \le d^{32}.
 \label{eq:binomial-integer122}
\end{equation}
If it holds, replace \(l\) by \(v\); otherwise replace \(h\) by \(v\).
The invariant \(B_k(l)\le\delta<B_k(h)\) holds initially and after every
step. At completion \(h-l=2^{-48}\) and
\(l\le L_k<h\). The reported endpoint is \(l\), an outward lower
bound, rather than a floating-point approximation to the root. For
\(k=0\) the reported lower endpoint is zero.

For each safe-improvement method and each of the two fixed thresholds
\(t\in\{2^{-24},2^{-20}\}\), set
\[
 I_{mj}(t)=\mathbf 1\{c_m(X_j)\ge t\},\qquad
 K_m(t)=\sum_{j=1}^{32} I_{mj}(t),\qquad
 p_m(t)=P_{40}\{c_m(X)\ge t\}.
\]
These are independent Bernoulli observations across state units,
although the 232 event series may be mutually dependent.
Equation~\eqref{eq:inference-safe-gain122} gives
\[
 \mathbb E_{40}G_m(X)\ge t\,p_m(t).
\]
Lemma~\ref{lem:binomial-inversion122} and the outward bisection
therefore justify \(t\,l_{K_m(t)}\) as a one-sided lower bound for
the original mean policy gain. An adaptive numerical refinement within
one state does not create another Bernoulli observation.

\begin{proposition}[Joint 95 percent conditional coverage]
\label{prop:joint-inference122}
Under the population and uniform numerical premises above, define
\[
 r_G=\frac{13\tau}{32},\qquad r_H=\frac{13\tau}{16}.
\]
For each of the 116 safe-improvement methods, use
\begin{align}
 \underline\mu_m
 &=\max\left\{0,\overline\ell_m-r_G,
                2^{-24}l_{K_m(2^{-24})},
                2^{-20}l_{K_m(2^{-20})}\right\},\\
 \overline\mu_m&=\min\{\tau,\overline u_m+r_G\}.
 \label{eq:joint-safe-interval122}
\end{align}
For each of the 348 prespecified signed comparisons, with paired sample
endpoints \(\ell^H_{mj},u^H_{mj}\), use
\begin{equation}
 [\,\max\{-\tau,\overline\ell^H_m-r_H\},
       \min\{\tau,\overline u^H_m+r_H\}\,],
 \label{eq:joint-signed-interval122}
\end{equation}
with \([0,0]\) for an exactly identical method key. With probability
at least \(19/20\), simultaneously all 116 intervals in
\eqref{eq:joint-safe-interval122}, all 232 event-probability lower
bounds, and all 348 signed intervals cover their stated population
quantities. The probability is conditional on the fixed
pre-acquisition information and is over the 32 independent initial
states, with no additional state units supplied by repeated clocks.
\end{proposition}

\begin{proof}
The normalized exponents in Lemma~\ref{lem:bounded-mean122} agree:
\[
 2(32)\left(\frac{r_G}{\tau}\right)^2
 =2(32)\left(\frac{r_H}{2\tau}\right)^2
 =\frac{169}{16}=:z.
\]
No numerical logarithm is needed to certify the family allocations.
The positive Taylor terms give
\[
 e^z>\sum_{j=0}^{30}\frac{z^j}{j!}=:T_{30}.
\]
For an independently checkable rational witness, the integer floors of
the first twelve terms, in increasing order of \(j\), are
\[
 1,\ 10,\ 55,\ 196,\ 518,\ 1095,\ 1928,\ 2910,\
 3842,\ 4509,\ 4763,\ 4573.
\]
Their sum is \(24400>23200\). Hence \(T_{30}>23200\), with all
comparisons obtainable by integer arithmetic. The union bound for the
116 two-sided safe mean bounds is at most
\[
 116\cdot2e^{-z}<232/T_{30}<1/100.
\]
For the 348 two-sided signed bounds it is at most
\[
 348\cdot2e^{-z}<696/T_{30}<3/100.
\]
The event inversions have total failure probability at most
\(232/23200=1/100\). The sum of these three family bounds is
\(1/20\). No independence between methods, comparisons, or event
thresholds is used.

On the joint event, each safe mean exceeds both its bounded-mean lower
endpoint and both of its threshold-derived lower endpoints, and is
below its bounded-mean upper endpoint. Their intersection therefore
retains coverage. The signed intervals follow directly from
\eqref{eq:inference-signed-support122} and
Lemma~\ref{lem:bounded-mean122}; an identical method has zero difference
at every state. This proves the proposition.
\end{proof}

\paragraph{Meaning and limits of the statement.}
The family covers fixed original-policy costs, with numerical endpoint
uncertainty retained. It neither proves a stronger claim than the
uniform qualification premises nor promotes empirical compatibility
to independent confirmation. The inherited actors are conditioned on;
no result here concerns their retraining on an unseen sample or their
optimization error under a different Bellman objective.

The complete-menu price domains use saved paired original-policy
intervals and recorded resource costs. They are descriptive conditional
scenarios, not additional confidence statements supplied by this
proposition. In particular, the timed scenarios use the fixed first
16 states for their economic means, whereas the simultaneous inference
above uses all 32 states. The three clocks, any restrictive additional
complete-clock sensitivity, reuse-volume scenarios, and repeated
administrative or statistical readers add zero independent population
observations. An own-incumbent improvement also does not identify a
neural-specific advantage over a traditional representation with the
same information, features, feasible policies, and resources.
